% chktex-file 12
% chktex-file 18
% chktex-file 36

\documentclass[12pt,twoside]{article}

\usepackage[a4paper,margin=1in]{geometry}
\usepackage{setspace}
\usepackage{fontspec}
\setmainfont{Times New Roman}
\usepackage{enumitem}
\usepackage[hidelinks]{hyperref}

% single spaced, blank line between paragraphs
\setlength{\parindent}{0pt}
\setlength{\parskip}{\baselineskip}
\setstretch{1}

% ---- APA 7 references with biblatex ----
\usepackage[
  backend=biber,
  style=apa,
  sorting=nyt
]{biblatex}
\DeclareLanguageMapping{english}{english-apa}
\addbibresource{references.bib}

\defbibheading{normalrefs}[\bibname]{%
  \setlength{\parskip}{0pt}% no extra blank line after heading
  \textbf{#1}\par
}

\usepackage{xcolor}
\usepackage{fancyhdr}
\pagestyle{fancy}
\fancyhf{} % clear all header/footer fields

% Optional: no header rule line
\renewcommand{\headrulewidth}{0pt}

% Make sure header has enough height for two lines
\setlength{\headheight}{26pt}

% Odd-page header, top left, gray text
\fancyhead[LO]{\small\textcolor{gray}{CTAW (Zhang)\\2025 Fall}}

% Footers: [page number] | [Title]. Odd pages: bottom left; Even pages: bottom right
\fancyfoot[LO]{\footnotesize\textcolor{gray}{\thepage\ \textbar\ “Us” vs “Them” in 1080p: Partisan Meme Wars on TikTok and X During the 2024 U.S. Election}}
\fancyfoot[RE]{\footnotesize\textcolor{gray}{\thepage\ \textbar\ “Us” vs “Them” in 1080p: Partisan Meme Wars on During the 2024 U.S. Election}}


\begin{document}

{\fontsize{16}{35}\selectfont
Critical Thinking \& Academic Writing\\
Project paper\par
}

\vspace{-0.5\baselineskip}

Class: C5 \\
Name: Yizhan Feng \\
Draft: D1 (2025.11.19) \\

\vspace{-0.75\baselineskip}

\begin{center}
{\fontsize{14}{17}\selectfont
``Us'' vs ``Them'' in 1080p: Partisan Meme Wars During the 2024 U.S. Election\par
}
\end{center}

\vspace{-\baselineskip}

% 1. Introduction

For many young Americans, the 2024 U.S. presidential election was experienced less through televised debates or newspaper columns and more through short videos and images on TikTok and X. Surveys by the Pew Research Center show that by 2024, about one-third of adults aged 18 to 29 regularly get their news from TikTok, with significant shares also using X and Instagram for political information. \parencite{shearer2024newsTikTokXFacebookInstagram} Within these platforms, memes---humorous images, short videos, and text formats that are rapidly remixed and shared---became a central medium for expressing political views. Media outlets compiled lists of ``the memes that defined the 2024 election season'', highlighting viral jokes about debate performances, aging candidates, migrant panics, and vice-presidential soundbites that turned into TikTok audio trends. 

However, many of these memes were not neutral jokes; they were partisan weapons, crafted and spread by supporters of each camp to hype their own side and mock the other. How do partisan memes on TikTok and X construct ``us'' and ``them'' in the 2024 U.S. election, and how is engagement with partisan meme warfare related to young users' sense of partisan identity and their willingness to engage with people from the opposing party? In light of these questions, this paper argues that partisan memes frame electoral politics as an antagonistic relationship between a virtuous, relatable ``us'' and a corrupt or ridiculous “them”. Among young users, heavy engagement with partisan meme content is associated with stronger partisan identity and lower openness to cross-party interaction, thereby contributing to affective polarization.

% 2. Literature Review

Although Internet users have only widely encountered ``meme'' in recent decades, the term itself was conceived long before the digital era. Derived from the Greek \textit{mimema}, the now famous term was first coined by \textcite{dawkins1976selfishgene} in his book \citetitle{dawkins1976selfishgene}, referring to small cultural units of transmission that, in analogy to genes, pass from person to person through various means of imitation. This definition was never static, however. The revival of meme in the Web 2.0 era further complicated it by attributing Internet memes with two fundamental characteristics: creative production and intertextuality \parencite{shifman2013memesdigitalculture}. The former refers to the adaptation and transformation of a meme through imitation or parody, either by mimicry or remix \parencite[20]{shifman2013memesdigitalculture}, while the latter emphasizes
how memes blend different cultural references or contexts. \textcite{shifman2013memesdigitalworld} further explored a three-dimensional typological framework for memetic analysis that breaks memes down into ``content, form, and stance'', emphasizing the power of stance in policing what emotions and positions are ``allowed'' and drawing boundaries between ``us'' and ``them''. 
 
% 3. Memetic Weaponization

Against this backdrop, Internet memes have played an increasingly visible role in recent U.S. presidential elections. Studies of the 2016 campaign show that supporters used image macros and other elements of remix culture on platforms such as Facebook and Reddit to circulate grassroots political messages, normalize pro-candidate narratives, and attack opponents through irony and ridicule \parencite{hunting2020popularmediamemes, moodyramirez2019facebookmemegroups}. The alt-right appropriation of Pepe the Frog, for instance, turned a benign cartoon character into a coded symbol of a virtuous, transgressive in-group supposedly fighting corrupt liberal elites \parencite{glitsos2020pepememeDarwinianAbsurd}. By 2020 and especially 2024, meme logics were no longer confined to fringe communities: major campaigns actively leaned into viral formats, as the Biden-Harris team did when it embraced the ``Dark Brandon'' meme and laser-eyed images of the president across TikTok and X to signal ironic strength and insider humor to young supporters \parencite{robertson2024bidenmemelord,mielczarek2024darkbrandonenthymeme}. On the Republican side, Donald Trump's catchphrase ``fake news'' itself had been repeatedly used not only as a meme template but also as a form of ``weaponized iconoclastic multimodal propaganda'' \parencite{smith2019weaponizediconoclasm} to portray Trump himself as the only truthful voice and construct a moral divide between a betrayed ``people'' and a corrupt ``establishment'' \parencite{ross2018discursivedeflection,woodward2020fakenewsguide}. Another highly visible Republican meme object was Trump's August 2023 mug shot from Fulton County, as journalism \parencite{yang2023mugshottomugs} and early meme studies \parencite{meleshchenko2024multimodalmetaphtonymy} documented how the image was instantly remixed into countless reaction memes, then appeared on T-shirts, mugs, and donation links, and eventually was reframed as proof that Trump is a persecuted yet defiant champion of ``the people''.

These memes were not neutral artifacts that simply popped up from nowhere and went viral randomly. Rather, they were deliberately selected and circulated as ``discursive weapons'' to ``justify judgement, condemnation, and exclusion of others'' \parencite[495]{nissenbaum2017contestedculturalcapital}. \textcite{peters2021weaponizingmemes} proposed this concept of memetic weaponization as ``purposeful deployment of memetic imagery to disrupt, undermine, attack, resist, or reappropriate discursive positions about public affairs in and around the news''. By serving as identificatory markers of affiliation, these memes successfully hijacked netizens' attention for political propaganda \parencite{olsen2018weaponizedmemes} and advanced a politics of othering, signaling who counts as in-group and who does not.
Taken together, these examples show how, by the 2024 election, both parties had weaponized memes as low-cost, high-affect tools for drawing the line between a morally superior, relatable ``us'' and an incompetent or dangerous ``them''. 

% 4. ``Coconut tree'' Case Analysis

One emblematic dynamic in this memetic warfare is the ``coconut tree'' meme around Kamala Harris. The meme originated from a May 10, 2023 White House speech, where Harris quoted her mother: ``You think you just fell out of a coconut tree? \dots You exist in the context of all in which you live and what came before you''---using the anecdote to argue that young people are shaped by families, communities, and history rather than existing ``in a silo'' \parencite{harris2023swearinginremarks}. This memefied fragment of Harris's speech was initially circulated by conservative Republican-aligned accounts as evidence of her supposed awkwardness and ``word salad'' communication style. Republicans deployed this meme as a vehicle of mockery, presenting Harris as incoherent, incomprehensible and fundamentally incapable. This strategy was effective, as the dominant reading on X was that Harris was ``weird'' or ``blissfully insane'', with users pairing the audio with reaction images to ridicule her \parencite{carry2024coconuttreeguide}. However, as Biden's position weakened, young Democrats online rallied to Harris as a possible replacement, and the meme flipped: \textit{The Week} described how coconut and palm-tree emojis became a shorthand for backing Harris, especially on X bios and display names \parencite{coleman2024coconutmeme}. Supporters of Harris even started calling themselves ``coconut-pilled'', adapting right-wing ``red-pilled'' language to signal that they had enthusiastically converted to Harris fandom \parencite{roberts2024coconuttreeorganic}.

This example is a robust validation of \textcite{shifman2013memesdigitalculture}'s memetic framework.  While the meme's content and form remained largely unchanged,
a simple shift in stance---from ``she's cringe, unserious, not presidential'' to ``she's a little weird, but in a fun, relatable, mom-wisdom way—and she's \textit{our} candidate''---allowed the same coconut tree clip to perform radically different discursive functions. Republicans weaponized the meme as an out-group attack focused on mockery and delegitimization, labeling ``them'' (Harris and the Democrats) as unserious, bizarre, and unfit. Democrats, on the other hand, embraced this hostile meme and reframed the quote as an empowering in-group hype, drawing a clear boundary between the playful, self-aware, community-rooted ``us'' and the unrelatable ``them'' who fail to understand the context.

% 5. Affective Polarization

Taken together, these examples suggest that politainment, a portmanteau word composed of politics and entertainment proposed by \textcite{doerner2001politainment} more than twenty years ago, has entered a new era of intensified online meme warfare. While scholars have documented the positive effect of political memes in opening up new participatory gateways for political engagement from below \parencite{sreekumar2013onlinepoliticalmemes, ahmed2024breakingbarriersmemes, mushtaq2025betweenhumorharm}, others warn that such participatory gain comes with important affective side-effects, namely affective polarization. Following \textcite{iyengar2012affectnotideology,iyengar2015fearloathing,iyengar2019originsaffectivepolarization}, affective polarization is defined as the growing tendency for partisans to evaluate their own side warmly while feeling dislike, distrust, and even disgust toward supporters of the opposing party. Rather than being driven purely by ideological distance, this form of polarization flows from social identity dynamics: seeing oneself as a loyal Democrat or Republican becomes a powerful basis for in-group favoritism and out-group animus. 

The 2024 election meme wars are likely to intensify this tendency. Quantitative analyses of mainstream political memes across recent U.S. election cycles show that the most heavily shared memes are strongly partisan in both content and stance, portraying the out-party as stupid, hypocritical, or dangerous while casting the in-party as sensible and morally righteous \parencite{collier2024compositionmainstreammemes}. As \textcite{alafnan2025roleMemesPoliticalDiscourse} notes, political memes ``serve as instruments for reinforcing ideological divides''. Rather than showing political opponents as ordinary citizens with different priorities, memes most often depict them as laughable or threatening caricatures, making the “other side” feel less legitimate and less human. By providing short runtimes and small canvases that encourage binary moral narratives, meme formats cloak hostility in humor and bake ``us vs them'' logics into their very form, further reinforcing the echo chamber effect. \textcite{kim2023partisanmemes}'s network studies confirm that engaging with partisan memes tends to catalyze homophilous online networks, as users who like or share such memes increasingly interact with like-minded others and see less cross-cutting content, reinforcing an “us versus them” structure in their feeds. Experimental work suggests that exposure to out-group-targeting memes does not so much persuade people to change sides as it entrenches existing identities, increasing polarization especially among strong partisans \parencite{galipeau2023impactpoliticalmemes}. Combined with platform algorithms that reward emotionally intense, in-group-affirming content \parencite{ye2025auditingexposurebias}, these dynamics mean that the same meme practices which draw young voters into politics also socialize them into experiencing politics as a zero-sum struggle against a contemptible “them”---precisely the emotional pattern captured by the concept of affective polarization.

% 6.Conclusion

Overall, partisan meme wars are double-edged. They energize and mobilize young voters, provide an outlet for anxiety and creativity, and help people feel that they are part of something larger. At the same time, they flatten complex issues into jokes and enemies, rewarding those who shout the loudest and mock the hardest. For democratic discourse, this raises a difficult question: how can societies harness the participatory power of memes without letting politics devolve into endless digital trench warfare? Answering this question requires not only platform governance and campaign ethics but also everyday critical literacy. Recognizing memes as simultaneously invitations to belong and cues to distrust may encourage young citizens, when they scroll past yet another Trump mugshot meme or coconut-tree clip, to pause before liking and reposting it and to ask instead what kind of ``us'' and ``them'' the meme is inviting them to imagine.

\vspace{-2\baselineskip}

\begin{flushright}\textbf{(1745 words)}\end{flushright}

\vspace{-\baselineskip}

\textbf{Declaration of generative AI and AI-assisted technologies in the writing process.}\\
During the preparation of this work, I used ChatGPT in order to brainstorm topics, summarize literature and generate title.
After using this tool/service, I reviewed and edited the content as needed and take(s) full responsibility for the content of the submitted project paper.
[Your reflection]

\vspace{\baselineskip}

\printbibliography[heading=normalrefs,title={References}]

\end{document}